\documentclass[11pt]{article}
\usepackage[margin=1in]{geometry}
\usepackage{amsmath,amssymb,amsthm,hyperref,enumitem}
\hypersetup{colorlinks=true,urlcolor=blue,linkcolor=blue,citecolor=blue}
\newtheorem{theorem}{Theorem}
\newtheorem{lemma}{Lemma}
\newcommand{\N}{\mathbb N}
\newcommand{\Q}{\mathbb Q}
\title{Almost all continued-fraction strings have\\ no nontrivial frequency symmetries}
\author{Principia Math In Progress\\MathDB 371294}
\date{18 September 2026}
\begin{document}
\maketitle
\begin{abstract}
For each fixed length $n\ge4$, at most $O_n(N^{n-1})$ strings with digits
in $\{1,\ldots,N\}$ have a permutation, other than the string and its
reverse, with the same Gauss--Kuzmin frequency. This proves Conjectures
4 and 6 of Singh, Zhang, and Hildebrand, including the unordered,
distinct-digit formulation in MathDB. The proof reconstructs the order
of independent variables, up to reversal, from one homogeneous part
of the characteristic polynomial.
\end{abstract}

\section*{Task 1: the full fixed-length theorem}
For a string $a=(a_1,\ldots,a_n)$ of positive integers, let $P_{\rm GK}(a)$
be the Gauss--Kuzmin measure of its continued-fraction cylinder.
A symmetry is nontrivial if it gives a permuted string different from
both $a$ and its reverse. Let $E_n(N)$ count such strings in
$\{1,\ldots,N\}^n$, allowing repeated digits. Let $f(N,n)$ count
$n$-element subsets of $\{1,\ldots,N\}$ having at least one ordering
with a nontrivial symmetry, and put
$\delta(N,n)=f(N,n)/\binom Nn$ for $N\ge n$.

\begin{theorem}\label{thm:main}
For every fixed integer $n\ge4$ and every positive integer $N$,
\[
 E_n(N)\le (n!-2)(2n-1)N^{n-1}.
\]
Consequently, for $N\ge n$,
\[
 \delta(N,n)\le
 \frac{(n!-2)(2n-1)N^{n-1}}{\binom Nn}
 =O_n(N^{-1})\longrightarrow0.
\]
\end{theorem}
The second conclusion is the exact MathDB assertion and
\cite[Conjecture 6, p.~29]{SZH}. The first also proves the
source's original formulation, Conjecture 4 on p.~7.

\subsection*{1. The known characteristic polynomial}
Define the continuant by
\[
 K()=1,\qquad K(x_1)=x_1,\qquad
 K(x_1,\ldots,x_m)=x_mK(x_1,\ldots,x_{m-1})
                         +K(x_1,\ldots,x_{m-2})\quad(m\ge2).
\]
For a string of length $n\ge3$, write
\[
 \begin{gathered}
 Q=K(x_1,\ldots,x_n),\quad
 U=K(x_1,\ldots,x_{n-1}),\quad V=K(x_2,\ldots,x_n),\\
 \chi(x_1,\ldots,x_n)=(Q+U)(Q+V).
 \end{gathered}
\]
This is the characteristic number of \cite[Definition 5.1]{SZH},
regarded as a polynomial in independent variables.
Indeed, if $p/q=[0;a_1,\ldots,a_n]$ and $p'/q'$ is the previous
convergent, the matrix identity
\[
 \prod_{j=1}^n
 \begin{pmatrix}0&1\\1&a_j\end{pmatrix}
 =
 \begin{pmatrix}p'&p\\q'&q\end{pmatrix}
\]
gives $q=Q(a)$, $q'=U(a)$, and $p=V(a)$.
The cylinder endpoints are $p/q$ and $(p+p')/(q+q')$.
Integrating $1/((1+x)\log2)$ between them gives the established formula
\cite[Lemma 2.5]{SZH}
\begin{equation}\label{eq:measure}
 P_{\rm GK}(a)=
 \left|\log_2\left(1+\frac{(-1)^n}{\chi(a)}\right)\right|.
\end{equation}
Here $\chi(a)>1$. For each fixed parity the expression on the right
is a strictly decreasing function of $\chi(a)$, so
\begin{equation}\label{eq:equivalence}
 P_{\rm GK}(\sigma(a))=P_{\rm GK}(a)
 \quad\Longleftrightarrow\quad
 \chi(\sigma(a))=\chi(a).
\end{equation}
This equivalence is \cite[Lemma 5.2]{SZH}; it is not a new claim here.

\subsection*{2. Recovering the order from high-degree terms}
Let $M=x_1\cdots x_n$. If $[F]_d$ denotes the homogeneous part of
total degree $d$, the continuant recurrence gives
\begin{equation}\label{eq:continuantparts}
 [Q]_n=M,\qquad [Q]_{n-1}=0,\qquad
 [Q]_{n-2}=\sum_{i=1}^{n-1}\frac{M}{x_ix_{i+1}}.
\end{equation}
To see the last formula by induction, the first term of the recurrence
extends each adjacent pair in the first $n-1$ positions; the second
term supplies the last pair $\{n-1,n\}$.
More generally, all degrees occurring in a length-$m$ continuant
have the same parity as $m$. Also
$[U]_{n-1}=M/x_n$ and $[V]_{n-1}=M/x_1$.
Expand $\chi=Q^2+Q(U+V)+UV$. The middle term has no part of degree
$2n-2$, by the parity observation. Thus
\begin{equation}\label{eq:key}
 [\chi]_{2n-2}
 =2\sum_{i=1}^{n-1}\frac{M^2}{x_ix_{i+1}}
       +\frac{M^2}{x_1x_n}.
\end{equation}
Every displayed quotient is a monomial with nonnegative exponents.
For $n\ge3$, the endpoint pair $\{1,n\}$ is different from all
adjacent pairs. Therefore the coefficient of $M^2/(x_ix_j)$, $i<j$,
is $2$ exactly when the two labels are adjacent in the ordering,
$1$ for its endpoint pair, and $0$ otherwise.

\begin{lemma}\label{lem:permutations}
For independent variables and $n\ge3$, the polynomial identity
\[
 \chi(x_{\sigma(1)},\ldots,x_{\sigma(n)})
       =\chi(x_1,\ldots,x_n)
\]
holds only when $\sigma$ is the identity or the reversal.
\end{lemma}
\begin{proof}
By \eqref{eq:key}, polynomial equality forces the two orderings to
have exactly the same unordered adjacent pairs. These pairs are the
edges of a path on the $n$ distinct variable labels. Such a path has
exactly two traversals visiting its vertices in order: start at one
endpoint or at the other. The two orderings are consequently equal
or reversed. Conversely, reversal preserves $\chi$, either by
transposing the displayed matrix product or by continuant reversal
symmetry.
\end{proof}
In particular, for every other permutation,
\[
 F_\sigma(x)=\chi(x_{\sigma(1)},\ldots,x_{\sigma(n)})-\chi(x_1,\ldots,x_n)
\]
is a nonzero polynomial. Its degree is at most $2n-1$, since the
leading term $[\chi]_{2n}=M^2$ cancels.

\subsection*{3. Counting integer zeros}
We include the elementary polynomial grid bound to make the counting
step self-contained.
\begin{lemma}\label{lem:grid}
A nonzero polynomial in $m$ variables over $\Q$ of total degree at
most $d$ has at most $dN^{m-1}$ zeros in $\{1,\ldots,N\}^m$.
\end{lemma}
\begin{proof}
For one variable this is the usual root bound.
For $m>1$, write
\[
 F=\sum_{j=0}^r F_j(x_1,\ldots,x_{m-1})x_m^j,\qquad F_r\ne0.
\]
The total degree of $F_r$ is at most $d-r$.
By induction, $F_r$ vanishes at at most $(d-r)N^{m-2}$ points of
the first $m-1$ coordinates. Over these points allow all $N$ values
of $x_m$; over every other point there are at most $r$ roots.
The total is at most
\[
 (d-r)N^{m-1}+rN^{m-1}=dN^{m-1}.
\]
Nonzero constants contribute no zeros, as required for the degree-zero
case.
\end{proof}

\begin{proof}[Proof of Theorem \ref{thm:main}]
If $a$ has a nontrivial symmetry, some index permutation $\sigma$
other than the identity and reversal satisfies $F_\sigma(a)=0$,
by \eqref{eq:equivalence}. Lemmas \ref{lem:permutations} and
\ref{lem:grid}, followed by a union bound over the $n!-2$ choices,
give
\[
 E_n(N)\le\sum_{\sigma\ne{\rm id},{\rm rev}}
             \#\{a\in\{1,\ldots,N\}^n:F_\sigma(a)=0\}
       \le(n!-2)(2n-1)N^{n-1}.
\]
Repeated digits cause no problem: the zero sets may count extra
tuples giving trivial string equalities, which only enlarges this
upper bound.
Each subset counted by $f(N,n)$ has at least one ordered witness
counted by $E_n(N)$, and different subsets have disjoint witnesses.
Hence $f(N,n)\le E_n(N)$.
Finally $\binom Nn\sim N^n/n!$ for fixed $n$, proving the result.
\end{proof}

\paragraph{Scope.}
The proof covers every fixed length and the source's intended
symmetry condition. It does not establish the stronger length-four
prediction $\log f(N,4)/\log N\to1$ in Conjecture 5(ii)
(MathDB 371295). That claim receives no solution credit.

\section*{Four proposed follow-up tasks}
\begin{enumerate}[label=\textbf{Task \arabic*:},start=2,leftmargin=*]
\item Prove or refute the sharper estimate $f(N,4)=N^{1+o(1)}$.
The present general argument gives only $f(N,4)=O(N^3)$.
\item Classify the positive integer solutions with four distinct
coordinates of $\chi(\sigma(a))=\chi(a)$, modulo permutation and
reversal, and determine whether finitely many parameterized families
account for all such solutions.
\item Determine the sharp growth exponents of $f(N,n)$ for each
fixed $n\ge5$, and identify which permutation equations contribute
the largest families of positive integer points.
\item Determine growth rates $n=n(N)\to\infty$ for which
$\delta(N,n(N))\to0$. Improve the factorial dependence in the
union bound by exploiting overlap and symmetry among the equations.
\end{enumerate}

\paragraph{Attribution and validation.}
The characteristic number, frequency equivalence, and reversal
symmetry are credited to \cite{SZH} and the classical sources cited
there. The contribution is the reconstruction in \eqref{eq:key}
and its use to prove the all-length zero-density assertion.
Exact computations check the polynomial identities, finite permutation
orbits, and the source's numerical counts. They supplement the
general proof rather than prove it by finite enumeration.
This manuscript has not received independent mathematical review.
The separate audits document the source and bounded prior-claim checks.

\begin{thebibliography}{1}
\bibitem{SZH} Shreyas Singh, Zhuo Zhang, and AJ Hildebrand,
\emph{An elementary characterization of the Gauss--Kuzmin measure
in the theory of continued fractions}, Integers \textbf{25} (2025),
Article A89, 32 pp., published 17 September 2025.
\href{https://math.colgate.edu/~integers/z89/z89.pdf}{Final published article};
DOI \href{https://doi.org/10.5281/zenodo.17144494}{10.5281/zenodo.17144494}.
See Lemmas 2.5 and 5.2, and Conjectures 4 and 6.
\end{thebibliography}
\end{document}
